\rtmtight
\begin{tblr}{width=\textwidth,colspec={|Q[c,m,2.1cm]|X[l,m]|},hlines,vlines,hspan=minimal,rowsep=1.5pt,colsep=3pt,row{1}={bg=headgray,font=\bfseries}}
\SetCell{c,m}\hfil\logo\hfil & \textbf{POLITEKNIK NEGERI MALANG}\par\textbf{JURUSAN TEKNOLOGI INFORMASI}\par\textbf{PROGRAM STUDI: D4 TEKNIK INFORMATIKA} \\
\SetCell[c=2]{l} \textbf{RENCANA TUGAS MAHASISWA (RTM) DAN RUBRIK PENILAIAN} & \\
Nama Mata Kuliah & Praktikum Basis Data \\
Kode Mata Kuliah & RTI252007 \quad \textbf{sks / Jam perminggu:} 2 SKS/4 jam \quad \textbf{Semester:} 2 \\
Dosen Pengampu & \begin{enumerate}\item Candra Bella Vista, S.Kom., M.T.\item Elok Nur Hamdana, S.T., M.T.\item M. Hasyim Ratsanjani, S.Kom., M.Kom.\item Ridwan Rismanto, S.ST., M.Kom., Ph.D.\item Vit Zuraida, S.Kom., M.Kom.\end{enumerate} \\
\end{tblr}
\assessmentblock{Laporan Praktikum Basis Data}{Laporan praktikum terstruktur (Case Based Learning)}{SCPMK0203-01601, SCPMK0701-01601}{Mahasiswa menyusun laporan praktikum untuk setiap pertemuan (Mgg 1--3, 5--8, 10--13) yang mencakup perancangan ERD, pemetaan model relasional, normalisasi, implementasi DDL, pengelolaan data DML, dan query SELECT sesuai studi kasus.}{Mengikuti praktikum terbimbing, mengerjakan latihan mandiri, mengimplementasikan rancangan pada MySQL, mendokumentasikan hasil eksekusi, dan menyusun laporan praktikum setiap pertemuan.}{Lembar laporan praktikum, diagram ERD, model relasional, skrip SQL, dan bukti eksekusi pada DBMS MySQL.}{Ketepatan perancangan ERD dan model relasional & 10 & Komponen ERD (entity, atribut, relationship, kardinalitas, partisipan) lengkap dan model relasional hasil pemetaan konsisten serta memenuhi 3NF. & Komponen utama ERD benar, tetapi sebagian kardinalitas atau normalisasi belum sepenuhnya tepat. & ERD tidak sesuai kebutuhan data atau model relasional masih menyisakan anomali. \\ \\
Ketepatan perintah SQL DDL dan DML & 10 & Perintah CREATE/ALTER/DROP dan INSERT/UPDATE/DELETE lengkap, sintaks benar, dan berhasil dieksekusi sesuai rancangan studi kasus. & Sebagian besar perintah benar, tetapi ada kesalahan sintaks atau constraint yang kurang tepat. & Perintah SQL tidak dapat dieksekusi atau tidak sesuai rancangan basis data. \\ \\
Kualitas laporan dan dokumentasi hasil eksekusi & 5 & Laporan runtut, bukti eksekusi lengkap, skrip SQL terdokumentasi rapi, dan hasil query sesuai kebutuhan studi kasus. & Laporan tersedia tetapi dokumentasi eksekusi atau kerapian skrip masih kurang. & Laporan tidak terdokumentasi atau hasil eksekusi tidak tercantum. \\}

\assessmentblock{Kuis Praktikum Basis Data}{Kuis ujian praktikum (closed book)}{SCPMK0203-01601, SCPMK0701-01601}{Mahasiswa mengerjakan Kuis Praktikum 1 (Mgg 4, materi Mgg 1--3: konsep basis data dan perancangan ERD) dan Kuis Praktikum 2 (Mgg 14, materi Mgg 10--13: SQL DDL, DML, DQL, dan JOIN) secara mandiri dan jujur.}{Mengerjakan soal ujian praktikum pilihan ganda dan/atau isian secara individu dengan sifat closed book sesuai jadwal asesmen.}{Lembar jawaban kuis praktikum.}{Kuis 1: penguasaan konsep basis data dan perancangan ERD & 7.5 & Jawaban tepat sesuai kunci; konsep basis data relasional, komponen ERD, dan langkah perancangan dijelaskan benar pada konteks soal. & Sebagian besar jawaban benar, tetapi ada kekeliruan dalam penerapan konsep pada soal. & Banyak jawaban tidak sesuai kunci atau konsep dasar basis data belum dikuasai. \\ \\
Kuis 2: penguasaan SQL DDL, DML, dan DQL & 7.5 & Jawaban tepat sesuai kunci; sintaks dan penggunaan perintah SQL (CREATE, INSERT, UPDATE, DELETE, SELECT, JOIN) benar sesuai soal. & Sebagian besar jawaban benar, tetapi ada kesalahan sintaks atau pemilihan perintah yang tidak tepat. & Banyak jawaban salah atau perintah SQL yang dituliskan tidak sesuai kebutuhan soal. \\}

\assessmentblock{Case Method Studi Kasus Basis Data}{Case Method dan demonstrasi unjuk kerja}{SCPMK0203-01601, SCPMK0701-01601, SCPMK1001-01601}{Mahasiswa menyelesaikan CM1 (Mgg 9, perancangan basis data dari studi kasus Mgg 1--8) dan CM2 (Mgg 15--17, studi kasus terpadu: identifikasi kebutuhan, perancangan ERD, pemetaan, normalisasi, implementasi MySQL, dan demonstrasi hasil).}{Menganalisis studi kasus, merancang ERD dan model relasional, mengimplementasikan pada MySQL, menyusun query DDL/DML/SELECT, mendokumentasikan laporan, dan mendemonstrasikan hasil kepada dosen.}{Laporan studi kasus, file ERD dan model relasional, skrip SQL, bukti eksekusi, dan demonstrasi langsung pada DBMS MySQL.}{Ketepatan identifikasi kebutuhan dan rancangan basis data & 20 & Kebutuhan data teridentifikasi lengkap, ERD konsisten dengan studi kasus, pemetaan benar, dan rancangan mencapai 3NF. & Kebutuhan utama teridentifikasi, tetapi sebagian rancangan atau bentuk normal belum sepenuhnya tepat. & Kebutuhan data tidak teridentifikasi dengan benar atau rancangan tidak sesuai studi kasus. \\ \\
Ketepatan implementasi basis data pada MySQL & 25 & Query DDL mengimplementasikan seluruh rancangan dengan benar termasuk relasi antartabel; DML mengisi dan memanipulasi data secara konsisten; SELECT (termasuk JOIN) menjawab kebutuhan informasi studi kasus. & Implementasi sebagian besar berhasil, tetapi ada tabel, relasi, atau query yang belum sesuai rancangan atau kebutuhan. & Implementasi gagal dieksekusi atau tidak sesuai hasil rancangan studi kasus. \\ \\
Kualitas demonstrasi dan laporan studi kasus & 15 & Demonstrasi berjalan lancar, mahasiswa mampu menjelaskan keputusan perancangan, dan laporan mendokumentasikan proses secara runtut dan lengkap. & Demonstrasi berhasil untuk fungsi utama, tetapi penjelasan keputusan atau kelengkapan laporan masih kurang. & Demonstrasi tidak berhasil atau laporan tidak menggambarkan proses pengerjaan studi kasus. \\}

\begin{tblr}{width=\textwidth,colspec={|Q[l,m,3.2cm]|X[l,m]|},hlines,vlines,hspan=minimal,row{1-3}={bg=headgray},rowsep=1.5pt,colsep=3pt}
Jadwal Pelaksanaan: & Minggu 1--17 sesuai RPS. Setiap evaluasi memiliki RTM dan rubrik multi-kriteria. \\
Lain-Lain Yang Diperlukan: & LMS, DBMS MySQL/MariaDB, perangkat pemodelan ERD, dan perangkat presentasi. \\
Pustaka: & Mengacu pustaka utama dan pendukung pada bagian RPS. \\
\end{tblr}
